%% LaTeX-Beamer template for KIT design
%% by Erik Burger, Christian Hammer
%% title picture by Klaus Krogmann
%%
%% version 2.1
%%
%% mostly compatible to KIT corporate design v2.0
%% http://intranet.kit.edu/gestaltungsrichtlinien.php
%%
%% Problems, bugs and comments to
%% burger@kit.edu

\documentclass[18pt]{beamer}

%% SLIDE FORMAT

% use 'beamerthemekit' for standard 4:3 ratio
% for widescreen slides (16:9), use 'beamerthemekitwide'

\usepackage{templates/beamerthemekit}
% \usepackage{templates/beamerthemekitwide}

%% TITLE PICTURE

% if a custom picture is to be used on the title page, copy it into the 'logos'
% directory, in the line below, replace 'mypicture' with the 
% filename (without extension) and uncomment the following line
% (picture proportions: 63 : 20 for standard, 169 : 40 for wide
% *.eps format if you use latex+dvips+ps2pdf, 
% *.jpg/*.png/*.pdf if you use pdflatex)

\titleimage{trellis}

%% TITLE LOGO

% for a custom logo on the front page, copy your file into the 'logos'
% directory, insert the filename in the line below and uncomment it

%\titlelogo{mylogo}

% (*.eps format if you use latex+dvips+ps2pdf,
% *.jpg/*.png/*.pdf if you use pdflatex)

%% TikZ INTEGRATION

% use these packages for PCM symbols and UML classes
% \usepackage{templates/tikzkit}
% \usepackage{templates/tikzuml}

% the presentation starts here

\title[A Generic Yet Fast Viterbi Decoder]{A Generic Yet Fast SIMD Viterbi-Decoder\\ Or My Try To Square The Circle}
\subtitle{}
\author{Jan Kraemer}

\institute{Communications Engineering Lab}

% Bibliography

\usepackage[citestyle=authoryear,bibstyle=numeric,hyperref,backend=biber]{biblatex}
\addbibresource{templates/example.bib}
\bibhang1em

\begin{document}

% change the following line to "ngerman" for German style date and logos
\selectlanguage{english}

%title page
\begin{frame}
\titlepage
\end{frame}

%table of contents
%\begin{frame}{Outline}
%\tableofcontents
%\end{frame}

\begin{frame}{Why speed up the decoding process}
\begin{itemize}
\item Viterbi Decoding is a very complex process
\begin{itemize}
	\item Neccessary for robust and fast communication
	\item Bottleneck regarding throughput
\end{itemize}
\pause
\item Gr-Trellis
\begin{itemize}
	\item Generic
	\item Not optimized
\end{itemize}
\pause
~ \newline
\item[$\Rightarrow$] More throughput would be nice!
\end{itemize}
\end{frame}

\begin{frame}{Gr-Trellis}
\begin{itemize}
\item Works with literally every trellis based encoding scheme  
\pause
\item Uses several lookup tables to be as generic as possible
\pause
\item Already implemented in GNURadio
\pause  
\item So let's just modify it!
\end{itemize}
\end{frame}


\begin{frame}{Amazing Results?}
\Huge NO!\\
~\\
\pause
\Large But Why?
\end{frame}
\begin{frame}{I blame the Lookup table}
$alpha[alphai*S+PS[j][i]]+in[k*O+OS[PS[j][i]*I+PI[j][i]]]$\\
~\\

\begin{itemize}
	\item SSE Operations work better with coherent memory\\~\\
	\item Generic lookup tables do not fulfill this requirement
\end{itemize}
\end{frame}


\begin{frame}{Start from scratch}
Assumptions\\~\\
	\begin{itemize}
		\item Rate 1/n codes
		\item Non-recursive codes
		\item Number of states $\geqslant$ 8	
	\end{itemize}
\end{frame}
\begin{frame}{Episode 4: A new decoder}
\begin{itemize}
	\item Integration into the new FEC-API	
	\item Based on Phillip Karns and FEC-APIs Voyager Decoder
	\item Uses a SSE float implementation 
\end{itemize}	
\end{frame}

\begin{frame}{The Add-Compare-Select Butterfly}

\begin{figure}[ht]
  \centering
  \includegraphics[scale=0.1]{logos/trellis_8.pdf}
\end{figure}

\end{frame}

\begin{frame}{The Add-Compare-Select Operation}

\begin{figure}[ht]
  \centering
  \includegraphics[scale=0.15]{logos/ACS.pdf}
\end{figure}

\end{frame}

\begin{frame}{Now for some real numbers}
\begin{figure}[]
  \includegraphics[scale=0.35]{logos/sse_vs_gen.pdf}
\end{figure}

\end{frame}

\begin{frame}{Conclusion}
\begin{itemize}
	\item SSE (float) is approx. 3.5x faster than standard C
	\item Memory constraints lead to reduced genericity
	\item Optimized Decoders are possible while a fair amount of flexibility is retained
	
\end{itemize}
\end{frame}

\begin{frame}{That's all folks}
	% minipage mit (Blind-)Text
	\begin{minipage}[t]{0.6\textwidth} 
	\begin{itemize}
	\item \textbf{Blog}:
	\begin{itemize}
	\item[]\textbf{spectrejan.blogspot.com}
	\end{itemize}
	\item \textbf{GITHUB}:
	\begin{itemize}
	\item[] \textbf{www.github.com/spectrejan}
	\end{itemize}
	\item \textbf{Email}:
	\begin{itemize}
	\item[] \textbf{kraemer.jn@gmail.com}
	\end{itemize}
	\end{itemize}
	\end{minipage}
	\begin{minipage}[t]{0.35\textwidth}
	\begin{minipage}[t]{1\textwidth}
	\begin{figure}[htbp]
	% \textwidth bezieht sich nun auf die Minipage
	\includegraphics[width=0.6\textwidth]{logos/GR}
	\end{figure} 
	\end{minipage}
	\begin{minipage}[t]{1\textwidth}
	\begin{figure}[htbp]
	% \textwidth bezieht sich nun auf die Minipage
	\includegraphics[width=0.6\textwidth]{logos/GSOC}
	\end{figure} 
	\end{minipage}
	\end{minipage}
\end{frame}

\end{document}
